% GENERATED FILE. Edit canonical lessons, then run npm run generate:books.
% canonical-source-hash: fnv1a64:46f8c9bca182c9b9
% canonical-lessons: JA-C243-ryoushin, JA-C243-sofu, JA-C243-sobo, JA-C243-kyoudai, JA-C243-sodateru

\chapter{Grandparents and Siblings}
\label{ch:ja-grandparents-and-siblings}
\hlchaptermodality{243}

\begin{chapteropening}
\textbf{By the end of this chapter:} \emph{I can ask and answer questions about my parents, grandparents and siblings.}

Complete the last lesson of chapter 243: I can ask and answer questions about my parents, grandparents and siblings.
\end{chapteropening}

\section[\texorpdfstring{\ja{りょうしん} (ryōshin)}{りょうしん (ryōshin)}]{\ja{りょうしん} (ryōshin) — parents; both parents}

\label{lesson:JA-C243-ryoushin}



\begin{quote}
Before the new one: say the Japanese for a close friend, then the Japanese for to go out with.
\end{quote}

\subsection*{You'll want to know: \ja{りょうしん}}
\textbf{\ja{りょうしん}} — \emph{ryōshin} — \textquotedblleft{}parents; both parents\textquotedblright{}.

\begin{quote}
\textbf{\ja{りょうしんは} \ja{げんきですか。}} — \emph{ryōshin wa genki desu ka} — \textquotedblleft{}Are our parents well?\textquotedblright{}
\end{quote}

\begin{cousinweb}
\textbf{\ja{りょうしん}} (\emph{ryōshin}) is Sino-Japanese (kango): two Chinese characters read with their on'yomi, readings borrowed from Chinese. They are \emph{ryō}, \textquotedblleft{}both\textquotedblright{} (Middle Chinese \emph{ljangX}) and \emph{shin}, \textquotedblleft{}parent, intimate\textquotedblright{} (\emph{tshin}) (KANJIDIC2; Middle Chinese in Baxter's transcription).
\end{cousinweb}

\begin{grammarlens}[title={What you've built}]
One more word for everyday life at A2.
\end{grammarlens}

\subsection*{Your turn}
\begin{itemize}
\raggedright
  \item \emph{Say it:} \emph{ryōshin}
  \item \emph{Say it:} \emph{ryōshin}, once more
  \item \emph{Say it:} \emph{ryōshin wa genki desu ka} — Are our parents well?
  \item \emph{Recall:} say the Japanese for a close friend, then the Japanese for to go out with, then say \emph{ryōshin} again
\end{itemize}

\begin{tcolorbox}[breakable,colback=teal!4,colframe=teal!35!black,title={Before you move on}]
What does \ja{りょうしん} mean? (\textquotedblleft{}Parents; both parents\textquotedblright{}.) Say it once more.
\end{tcolorbox}
\section[\texorpdfstring{\ja{そふ} (sofu)}{そふ (sofu)}]{\ja{そふ} (sofu) — my grandfather; a grandfather}

\label{lesson:JA-C243-sofu}



\begin{quote}
Before the new one: say the Japanese for to go out with, then the Japanese for parents.
\end{quote}

\subsection*{You'll want to know: \ja{そふ}}
\textbf{\ja{そふ}} — \emph{sofu} — \textquotedblleft{}my grandfather; a grandfather\textquotedblright{}.

\begin{quote}
\textbf{\ja{そふは} \ja{もう} \ja{はたらきません。}} — \emph{sofu wa mō hatarakimasen} — \textquotedblleft{}My grandfather no longer works.\textquotedblright{}
\end{quote}

\begin{cousinweb}
\textbf{\ja{そふ}} (\emph{sofu}) is Sino-Japanese (kango): two Chinese characters read with their on'yomi, readings borrowed from Chinese. They are \emph{so}, \textquotedblleft{}ancestor\textquotedblright{} (Middle Chinese \emph{tsuX}) and \emph{fu}, \textquotedblleft{}father\textquotedblright{} (\emph{bjuX}) (KANJIDIC2; Middle Chinese in Baxter's transcription).
\end{cousinweb}

\begin{grammarlens}[title={What you've built}]
One more word for everyday life at A2.
\end{grammarlens}

\subsection*{Your turn}
\begin{itemize}
\raggedright
  \item \emph{Say it:} \emph{sofu}
  \item \emph{Say it:} \emph{sofu}, once more
  \item \emph{Say it:} \emph{sofu wa mō hatarakimasen} — My grandfather no longer works.
  \item \emph{Recall:} say the Japanese for to go out with, then the Japanese for parents, then say \emph{sofu} again
\end{itemize}

\begin{tcolorbox}[breakable,colback=teal!4,colframe=teal!35!black,title={Before you move on}]
What does \ja{そふ} mean? (\textquotedblleft{}My grandfather; a grandfather\textquotedblright{}.) Say it once more.
\end{tcolorbox}
\section[\texorpdfstring{\ja{そぼ} (sobo)}{そぼ (sobo)}]{\ja{そぼ} (sobo) — my grandmother; a grandmother}

\label{lesson:JA-C243-sobo}



\begin{quote}
Before the new one: say the Japanese for parents, then the Japanese for my grandfather.
\end{quote}

\subsection*{You'll want to know: \ja{そぼ}}
\textbf{\ja{そぼ}} — \emph{sobo} — \textquotedblleft{}my grandmother; a grandmother\textquotedblright{}.

\begin{quote}
\textbf{\ja{そぼは} \ja{にくを} \ja{たべません。}} — \emph{sobo wa niku o tabemasen} — \textquotedblleft{}My grandmother does not eat meat.\textquotedblright{}
\end{quote}

\begin{cousinweb}
\textbf{\ja{そぼ}} (\emph{sobo}) is Sino-Japanese (kango): two Chinese characters read with their on'yomi, readings borrowed from Chinese. They are \emph{so}, \textquotedblleft{}ancestor\textquotedblright{} (Middle Chinese \emph{tsuX}) and \emph{bo}, \textquotedblleft{}mother\textquotedblright{} (\emph{muwX}) (KANJIDIC2; Middle Chinese in Baxter's transcription).
\end{cousinweb}

\begin{grammarlens}[title={What you've built}]
One more word for everyday life at A2.
\end{grammarlens}

\subsection*{Your turn}
\begin{itemize}
\raggedright
  \item \emph{Say it:} \emph{sobo}
  \item \emph{Say it:} \emph{sobo}, once more
  \item \emph{Say it:} \emph{sobo wa niku o tabemasen} — My grandmother does not eat meat.
  \item \emph{Recall:} say the Japanese for parents, then the Japanese for my grandfather, then say \emph{sobo} again
\end{itemize}

\begin{tcolorbox}[breakable,colback=teal!4,colframe=teal!35!black,title={Before you move on}]
What does \ja{そぼ} mean? (\textquotedblleft{}My grandmother; a grandmother\textquotedblright{}.) Say it once more.
\end{tcolorbox}
\section[\texorpdfstring{\ja{きょうだい} (kyōdai)}{きょうだい (kyōdai)}]{\ja{きょうだい} (kyōdai) — siblings; brothers and sisters}

\label{lesson:JA-C243-kyoudai}



\begin{quote}
Before the new one: say the Japanese for my grandfather, then the Japanese for my grandmother.
\end{quote}

\subsection*{You'll want to know: \ja{きょうだい}}
\textbf{\ja{きょうだい}} — \emph{kyōdai} — \textquotedblleft{}siblings; brothers and sisters\textquotedblright{}.

\begin{quote}
\textbf{\ja{きょうだいが} \ja{いますか。}} — \emph{kyōdai ga imasu ka} — \textquotedblleft{}Do you have any brothers or sisters?\textquotedblright{}
\end{quote}

\begin{cousinweb}
\textbf{\ja{きょうだい}} (\emph{kyōdai}) is Sino-Japanese (kango): two Chinese characters read with their on'yomi, readings borrowed from Chinese. They are \emph{kyō}, \textquotedblleft{}elder brother\textquotedblright{} (Middle Chinese \emph{xjwaeng}) and \emph{dai}, \textquotedblleft{}younger brother\textquotedblright{} (\emph{dejX}) (KANJIDIC2; Middle Chinese in Baxter's transcription).
\end{cousinweb}

\begin{grammarlens}[title={What you've built}]
One more word for everyday life at A2.
\end{grammarlens}

\subsection*{Your turn}
\begin{itemize}
\raggedright
  \item \emph{Say it:} \emph{kyōdai}
  \item \emph{Say it:} \emph{kyōdai}, once more
  \item \emph{Say it:} \emph{kyōdai ga imasu ka} — Do you have any brothers or sisters?
  \item \emph{Recall:} say the Japanese for my grandfather, then the Japanese for my grandmother, then say \emph{kyōdai} again
\end{itemize}

\begin{tcolorbox}[breakable,colback=teal!4,colframe=teal!35!black,title={Before you move on}]
What does \ja{きょうだい} mean? (\textquotedblleft{}Siblings; brothers and sisters\textquotedblright{}.) Say it once more.
\end{tcolorbox}
\section[\texorpdfstring{\ja{そだてる} (sodateru)}{そだてる (sodateru)}]{\ja{そだてる} (sodateru) — to bring up; to raise}

\label{lesson:JA-C243-sodateru}



\begin{quote}
Before the new one: say the Japanese for my grandmother, then the Japanese for siblings.
\end{quote}

\subsection*{You'll want to know: \ja{そだてる}}
\textbf{\ja{そだてる}} — \emph{sodateru} — \textquotedblleft{}to bring up; to raise\textquotedblright{}.

\begin{quote}
\textbf{\ja{あねは} \ja{こどもを} \ja{ひとりで} \ja{そだてますか。}} — \emph{ane wa kodomo o hitori de sodatemasu ka} — \textquotedblleft{}Will my elder sister bring up her child on her own?\textquotedblright{}
\end{quote}

Polite \textbf{\ja{そだてます}}, \emph{sodatemasu}; an ichidan verb.

\begin{cousinweb}
\textbf{\ja{そだてる}} (\emph{sodateru}) is a native Japanese word (wago), inherited rather than borrowed from Chinese; its kanji was attached later as a kun'yomi (JMdict, KANJIDIC2).
\end{cousinweb}

\begin{grammarlens}[title={What you've built}]
One more word for everyday life at A2.
\end{grammarlens}

\subsection*{Your turn}
\begin{itemize}
\raggedright
  \item \emph{Say it:} \emph{sodateru}
  \item \emph{Say it:} \emph{sodateru}, once more
  \item \emph{Say it:} \emph{ane wa kodomo o hitori de sodatemasu ka} — Will my elder sister bring up her child on her own?
  \item \emph{Recall:} say the Japanese for my grandmother, then the Japanese for siblings, then say \emph{sodateru} again
\end{itemize}

\begin{tcolorbox}[breakable,colback=teal!4,colframe=teal!35!black,title={Before you move on}]
What does \ja{そだてる} mean? (\textquotedblleft{}To bring up; to raise\textquotedblright{}.) Say it once more.
\end{tcolorbox}
